\begin{frame}{BN의 한계: 작은 배치 — LayerNorm/InstanceNorm/GroupNorm}
  \small
  \begin{itemize}
    \item BN은 \textbf{batch 통계}에 의존 — batch size가 \textbf{1$\sim$2}
      로 작으면 통계가 극도로 불안정, BN이 제대로 동작 안 함
    \item \kb{LayerNorm}: 배치가 아니라 \textbf{특징(feature) 축}을
      정규화 — 시계열·NLP(트랜스포머, 작은 배치/변동 길이) 표준
    \item \kb{InstanceNorm}: batch를 무시, 샘플·채널 단위 (생성 모델)
    \item \kb{GroupNorm}: 채널을 그룹 단위로 묶어 정규화 —
      batch size에 의존하지 않음
  \end{itemize}
  \vskip0.3em
  \begin{block}{``BN을 쓰면 무조건 좋은가?''}
    \textbf{배치가 충분히 큰 분류/회귀 CNN}에서는 대체로 Yes,
    \textbf{그렇지 않은 곳}에서는 LayerNorm 계열.
  \end{block}
  \vskip0.2em
  {\footnotesize \kb{자주 하는 실수}: batch 1$\sim$2 (시계열 RNN,
  배치 1 추론)에서 BN을 고집 — batch 통계가 1$\sim$2개 샘플로
  추정돼 분산이 0 근처 폭주/무의미. 이 경우 \textbf{LayerNorm이
  정석}.}
\end{frame}
